\documentclass[11pt]{amsart}
\usepackage{calc,amssymb,amsmath,ulem,stmaryrd,fullpage}
\usepackage{enumerate}
\usepackage{alltt}
\RequirePackage[dvipsnames,usenames]{color}
\let\mffont=\sf
\normalem
%\input{kmacros3.sty}
\input{xy}
\xyoption{all}
\usepackage{hyperref}
\usepackage{graphicx}
\usepackage[all,cmtip]{xy}
\usepackage{verbatim}
\usepackage[left=0.95in,top=0.95in,right=0.95in,bottom=0.95in]{geometry}
\newcommand{\tauCohomology}{T}
\newcommand{\FCohomology}{S}


\theoremstyle{theorem}
%\newtheorem*{mainthm*}{Main Theorem}

\newcommand{\note}[1]{\marginpar{\sffamily\tiny #1}}

\def\todo#1{\textcolor{Mahogany}%
{\sffamily\footnotesize\newline{\color{Mahogany}\fbox{\parbox{\textwidth-15pt}{#1}}}\newline}}

\renewcommand{\O}{\mathcal O}
\newcommand{\ie}{{\itshape i.e.} }
\newcommand{\roundup}[1]{\lceil #1 \rceil}
\newcommand{\rounddown}[1]{\lfloor #1 \rfloor}
\renewcommand{\to}[1][]{\xrightarrow{\ #1\ }}
\newcommand{\onto}[1][]{\protect{\xrightarrow{\ #1\ }\hspace{-0.8em}\rightarrow}}
\newcommand{\into}[1][]{\lhook \joinrel \xrightarrow{\ #1\ }}
%\newcommand{DBDual}[1]{{\underline{\omega}_{#1}}}
\newcommand{\DBDual}[1]{{{\uuline \omega}{}^{\mydot}_{#1}}}
\newcommand{\Tt}{{\mathfrak{T}}}
%\newcommand{\bn}{{\mathfrak{n}} }
\newcommand{\sol}[1]{ \vskip 6pt{\color{blue} {\bf Solution:} #1} \vskip 6pt}

\newcommand{\bR}{{\mathbb{R}}}
\newcommand{\va}{{\vec{a}}}
\newcommand{\vb}{{\vec{b}}}
\newcommand{\vzero}{{\vec{0}}}
\newcommand{\vi}{{\vec{i}}}
\newcommand{\vj}{{\vec{j}}}
\newcommand{\vk}{{\vec{k}}}
\newcommand{\vh}{{\vec{h}}}
\newcommand{\vr}{{\vec{r}}}
\newcommand{\vu}{{\vec{u}}}
\newcommand{\vv}{{\vec{v}}}
\newcommand{\vw}{{\vec{w}}}
\newcommand{\vx}{{\vec{x}}}
\newcommand{\vy}{{\vec{y}}}
\newcommand{\vz}{{\vec{z}}}
\newcommand{\vT}{{\vec{T}}}
\newcommand{\vN}{{\vec{N}}}
\newcommand{\vF}{{\vec{F}}}
\newcommand{\vG}{{\vec{G}}}
\newcommand{\bF}{\mathbb{F}}
\newcommand{\bQ}{\mathbb{Q}}
\newcommand{\bZ}{\mathbb{Z}}
\newcommand{\Inn}{\mathrm{Inn}}
\newcommand{\Aut}{\mathrm{Aut}}


\begin{document}
\title{Worksheet \#2 -- Math 6310\\Fall 2017}
\author{Due Wednesday October 18th}
\maketitle

We'll have some fun with localization.
\vskip 9pt
\noindent
{\bf Definition 1.}  Suppose $R$ is a integral domain.  A subset $W \subseteq R$ containing $1$ and not containing $0$ is called a \emph{multiplicative set} if for any $a,b \in W$, we have $ab \in W$.  (Note the conditions on 1 and 0 are not always assumed in the literature, we make them here for convenience).
\vskip 9pt
\noindent
{\bf 1.}  Suppose that $Q \subseteq R$ is a prime ideal.  Show that $R \setminus Q$ is a multiplicative set.  In particular, note that $R \setminus \{0 \}$ is a multiplicative set in an integral domain.
\vskip 150pt
\noindent
{\bf Definition 2.}  Suppose $W \subseteq R$ is a multiplicative set in an integral domain.  Consider the set of all pairs $\{(r, w) \;|\; r \in R, w \in W \}$ under the following equivalence relation $(r, w) \sim (r', w')$ if $w' r = w r'$.  The set of equivalence classes is denoted $W^{-1} R$ and equivalence classes of pairs $(r, w)$ are denoted by $r/w$.
\vskip 9pt
\noindent
{\bf 2.}  Suppose that $W \subseteq R$ is a multiplicative set.  Show that $W^{-1} R$ is a ring under the operations induced by the formulae
\[
r/w + r'/w' = (r'w + rw')/(ww') \text{ and } (r/w)(r'/w') = (rr'/ww').
\]
In particular, show that those operations are well defined.  Also show that the map $\phi : R \to W^{-1} R$ defined by $\phi(r) = r/1$ is an injective ring homomorphism.
\newpage
\noindent
{\bf 3.}  If $R$ is an integral domain and $W = R \setminus \{ 0 \}$ show that $W^{-1}R$ is a field.  This field is called the field of fractions of $R$.
\vskip 100pt
\noindent
{\bf 4.}  If $R$ is an integral domain, $W = R \setminus \{ 0 \}$, and $K = W^{-1} R$ is the field of fractions, show that $K$ contains an isomorphic copy of $R$ as a subring.
\vskip 100pt
\noindent
{\bf 5.}  Now suppose that $R$ is an integral domain, $W \subseteq R$ is a multiplicative set and $Q \subseteq R$ is a prime ideal such that $W \cap Q = \emptyset$.  Show that the set
\[
W^{-1} Q := \{ x/w \;|\; x \in Q, w \in W \} \subseteq W^{-1} R
\]
is a prime ideal of $W^{-1} R$.
\newpage
\noindent
{\bf 6.}  Suppose that $R$ is an integral domain and $W \subseteq R$ is a multiplicative set.  Suppose that $P \subseteq W^{-1} R$ is a prime ideal and set $Q = P \cap R$ (where the intersection is taken with the isomorphic copy of $R \cong \phi(R)$ from {\bf 2.}).  Show that $Q \cap W = \emptyset$ and that $W^{-1} Q = P$.
\vskip 200pt
\noindent
{\bf 7.}  Conclude that the prime ideals of $W^{-1} R$ are in bijection with the prime ideals $Q$ in $R$ such that $W \cap Q = \emptyset$.
\vskip 150pt
\noindent
{\bf 8.}  Suppose that $W \subseteq R$ is a multiplicative set in an integral domain made up of only units (= invertible elements) of $R$.  Show that the map $\phi : R \to W^{-1} R$ you constructed in {\bf 2.} is an isomorphism.
\newpage
\noindent
{\bf 9.}  Give an example of an integral domain and two multiplicative sets $W, V \subseteq R$ such that $W \neq V$, $W^{-1} R \cong V^{-1} R$ but $\phi : R \to W^{-1} R$ is not surjective.
\vskip 200pt
\noindent
{\bf 10.}  Suppose that $R$ is an integral domain, $W \subseteq R$ is a multiplicative set and $\psi : R \to S$ is a ring homomorphism between commutative rings.  Further suppose that for each $w \in W$, there is some $s \in S$ such that $s \psi(w) = 1_S$.  Show that one can factor $\psi$ as follows
\[
\xymatrix{
R \ar[dr]_-{\phi} \ar[rr]^{\psi} & & S\\
& W^{-1} R \ar@{.>}[ur]_-{\alpha}
}
\]
and that the map $\alpha$ is unique.  This is called the \emph{universal property of localization}.
\end{document}

